\corrigesChapitre{7}{Intégrales, aires et valeur moyenne}
\begin{corrige}{1}{Calculer avec une primitive}
\begin{questions}
\item $[x^3]_0^2=8$.
\item $[\ln x]_1^{\mathrm e}=1$.
\item $[\mathrm e^{2x}/2]_0^1=(\mathrm e^2-1)/2$.
\end{questions}
\end{corrige}
\begin{corrige}{2}{Chasles et linéarité}
\begin{questions}
\item $3+7=10$.
\item $2\times10+(5-0)=25$.
\item $-10$, car inverser les bornes change le signe.
\end{questions}
\end{corrige}
\begin{corrige}{3}{Valeur moyenne}
\begin{questions}
\item $\frac13[x^3/3]_0^3=3$.
\item La moyenne des bornes vaut $(0+9)/2=4{,}5$, différente.
\item Oui : $0\leq3\leq9$.
\end{questions}
\end{corrige}
\begin{corrige}{4}{Une aire signée}
\begin{questions}
\item $[x^2/2-2x]_0^4=0$.
\item Deux triangles de base $2$ et hauteur $2$, donc $2+2=4$.
\item L’intégrale retranche la zone négative ; les deux zones se compensent. L’aire totale les additionne.
\end{questions}
\end{corrige}
\begin{corrige}{5}{Entre deux courbes}
\begin{questions}
\item $x-x^2=x(1-x)\geq0$, donc $f\geq g$.
\item $\int_0^1(x-x^2)\,\mathrm dx=1/2-1/3=1/6$.
\item $x=x^2$ donne $x=0$ et $x=1$, soit $(0,0)$ et $(1,1)$.
\end{questions}
\end{corrige}
\begin{corrige}{6}{Production cumulée}
\begin{questions}
\item $[-12\mathrm e^{-t/2}]_0^4=12(1-\mathrm e^{-2})\simeq10{,}376$ litres.
\item $3(1-\mathrm e^{-2})\simeq2{,}594$ L/min.
\item Le débit décroît de $6$ à $6\mathrm e^{-2}\simeq0{,}812$ ; la moyenne est bien entre ces bornes.
\end{questions}
\end{corrige}
\begin{corrige}{7}{Rectangles}
\begin{questions}
\item $\frac14(0+1/16+1/4+9/16)=7/32=0{,}21875$.
\item À droite : $15/32=0{,}46875$ ; trapèzes : $11/32=0{,}34375$.
\item L’exact vaut $1/3$. Les rectangles l’encadrent ; les trapèzes le surestiment ici de $1/96$.
\end{questions}
\end{corrige}
\begin{corrige}{8}{Lorenz et Gini}
\begin{questions}
\item $L(1/2)=1/8=12{,}5\,\%$.
\item $G=1-2[x^4/4]_0^1=1/2$.
\item Pour $x^2$, $G=1/3$. Le modèle $x^3$ présente ici davantage d’inégalités, ce qui concorde avec l’ordre des courbes.
\end{questions}
\end{corrige}
\begin{corrige}{9}{Une estimation aléatoire}
\begin{questions}
\item $3\,340/10\,000=0{,}334$.
\item L’exact est $1/3\simeq0{,}333333$ ; l’erreur absolue observée est environ $0{,}000667$.
\item Non. Une simulation fournit une estimation fluctuante ; la preuve exacte utilise ici une primitive.
\end{questions}
\end{corrige}
\begin{corrige}{10}{Données groupées}
\begin{questions}
\item Premier trapèze : $0{,}5(0+0{,}2)/2=0{,}05$ ; second : $0{,}5(0{,}2+1)/2=0{,}30$. Total $0{,}35$.
\item $G=1-2(0{,}35)=0{,}30$.
\item Les parts cumulées de deux groupes ne déterminent pas la répartition interne à chaque groupe ; la ligne polygonale est une représentation groupée.
\end{questions}
\end{corrige}
\begin{corrige}{11}{Primitive définie par une aire}
\begin{questions}
\item $H(1)=0$ et $H\,\!^{\prime}(x)=x^2+1$ par le théorème fondamental.
\item $H(x)=x^3/3+x-4/3$.
\item La dérivée est strictement positive sur $\R$ ; $H$ est strictement croissante.
\end{questions}
\end{corrige}
\begin{corrige}{12}{Encadrer sans primitive}
\begin{questions}
\item La dérivée $-2x\mathrm e^{-x^2}$ est négative ou nulle sur $[0,1]$ ; la fonction décroît.
\item Les rectangles à droite minorent ; ceux à gauche majorent.
\item La différence télescopique vaut $(1-\mathrm e^{-1})/n$. Il suffit de $n\geq633$, car $1-\mathrm e^{-1}\simeq0{,}632121$.
\end{questions}
\end{corrige}
